\section{Output Analysis}
\label{sec:output}

This is a \textbf{terminating} simulation: the system starts empty-and-idle and
ends when the 50\textsuperscript{th} batch ships (or is scrapped), so no warm-up
period is deleted. Each scenario is run as $n$ independent replications with
unique random seeds in a \texttt{Parameter Variation} experiment.

\paragraph{Number of replications.}
We size $n$ with the half-width method on the headline KPI (weekly throughput in
bottles shipped). From $n_0 = 10$ pilot replications we apply Rossetti's formula:
\begin{equation}
  n \;\ge\; \left(\frac{t_{n_0-1,\;1-\alpha/2}\;s_0}{\varepsilon}\right)^{\!2},
  \label{eq:halfwidth}
\end{equation}
targeting relative precision $\varepsilon \le \SI{5}{\percent}$ of the mean at
$\alpha = 0.05$. Table~\ref{tab:reps} shows the sizing results from the
full $N=20$ baseline run.

\begin{table}[H]
\centering
\caption{Half-width analysis from $N=20$ baseline replications (95\,\% CI,
         $t_{19,\,0.975}=2.093$). Scrap rate uses an absolute target
         $\varepsilon_\text{abs}=0.03$ because the relative target is
         unattainable within a feasible run count.}
\label{tab:reps}
\small
\begin{tabularx}{\textwidth}{X r r r r r}
\toprule
KPI & Mean & Std dev & Half-width & Rel.\ h-w & $n$ required \\
\midrule
Bottles shipped       & 67\,487  & 2\,725   & $\pm$1\,275  & 1.9\,\%  & 3   \\
Batches shipped       & 46.8     & 1.74     & $\pm$0.81    & 1.7\,\%  & 3   \\
Mean flow time (h)    & 200.9    & 7.95     & $\pm$3.72    & 1.9\,\%  & 3   \\
Max flow time (h)     & 376.8    & 15.7     & $\pm$7.37    & 2.0\,\%  & 4   \\
Avg WIP (batches)     & 19.85    & 0.73     & $\pm$0.34    & 1.7\,\%  & 3   \\
Dryer utilisation     & 0.946    & 0.008    & $\pm$0.004   & 0.4\,\%  & 1   \\
Cleaning-team util.   & 0.819    & 0.013    & $\pm$0.006   & 0.7\,\%  & 1   \\
Packaging util.       & 0.545    & 0.006    & $\pm$0.003   & 0.5\,\%  & 1   \\
Scrap rate            & 0.060    & 0.035    & $\pm$0.016   & 25.4\,\% & 516 \\
\bottomrule
\end{tabularx}
\end{table}

All KPIs except scrap rate require $n \le 5$ for 5\,\% relative precision;
scrap rate is highly variable because a single replication can scrap up to 11
of 50 batches. Requiring 5\,\% relative precision on scrap would demand
$n \approx 516$ replications, which is infeasible. We therefore adopt
$N = 20$ as our final run size: this yields tight confidence intervals on
throughput, flow time, and utilisation ($\le 2.1\,\%$ relative half-width),
and an absolute half-width of $\pm 0.016$ on scrap rate, sufficient for
scenario comparisons. All scenario runs use the same $N = 20$.
